\chapter{Wnioski i kierunki dalszych badań}
\label{chap:wnioski}

\section{Podsumowanie wyników}

Celem niniejszej pracy była analiza narzędzi i metod wykrywania treści generowanych
przez sztuczną inteligencję w obszarze multimediów wideo, ze szczególnym uwzględnieniem
praktycznej skuteczności dostępnych detektorów oraz zaproponowanie podejścia
wielosygnałowego jako alternatywy dla rozwiązań komercyjnych.

Przeprowadzone eksperymenty wykazały, że popularne narzędzie DetectVideo.ai zwraca
wynik 55\% (inconclu\-sive) dla ponad 80\% testowanych materiałów, w tym dla nagrań
autentycznych, co czyni je praktycznie bezużytecznym do binarnej klasyfikacji treści.
Jedynie filmy o wyraźnych artefaktach wizualnych typowych dla generatywnych
modeli (np.\ abstrakcyjne symulacje płynów, sceny fantasy) uzyskały wynik 72\%,
co nadal nie przekracza przyjętego w literaturze progu pewności detekcji.
Proponowane podejście wielosygnałowe --- łączące analizę spektralną, przepływ
optyczny oraz analizę metadanych --- wykazuje zdolność do podjęcia jednoznacznej
decyzji binarnej dla każdego badanego pliku, co stanowi jego kluczową przewagę
praktyczną.

\section{Kierunki dalszych badań}

Niniejsza praca stanowi punkt wyjścia do szeroko zakrojonych badań w obszarze
detekcji multimediów generowanych przez AI. Poniżej przedstawiono kierunki,
które zdaniem autora posiadają największy potencjał naukowy i praktyczny.

\subsection{Rozszerzenie zbioru danych testowych}

Podstawowym ograniczeniem przeprowadzonych eksperymentów jest stosunkowo niewielki
zbiór 24 plików wideo. Dalsze badania powinny obejmować:

\begin{itemize}
  \item \textbf{Większe zbiory darmowe} --- takie jak FaceForensics++~\cite{Rossler2019FaceForensics},
        DFDC (Deepfake Detection Challenge Dataset) czy WildDeepfake, liczące dziesiątki
        tysięcy klipów z różnymi technikami syntezy;
  \item \textbf{Płatne bazy danych} --- komercyjne zbiory, np.\ oferowane przez Veritone
        lub Nielsen, zawierające materiały z kontrolowanymi metadanymi produkcji,
        co umożliwia precyzyjną weryfikację etykiet;
  \item \textbf{Materiały z platform społecznościowych} --- filmy pobrane z TikToka,
        YouTube Shorts i Instagrama, które przeszły przez agresywną rekompresję
        transkodowania platformy, stanowiącą jedno z najtrudniejszych wyzwań
        dla detektorów sygnałowych;
  \item \textbf{Zróżnicowanie modeli generatywnych} --- włączenie materiałów z nowszych
        modeli (Sora 2, Runway Gen-5, Kling, Pika 2.0) w miarę ich dostępności,
        aby śledzić ewolucję artefaktów.
\end{itemize}

\subsection{Ewaluacja płatnych narzędzi komercyjnych}

W ramach niniejszej pracy przetestowano wyłącznie narzędzia dostępne bezpłatnie.
Pełna analiza porównawcza powinna objąć płatne rozwiązania klasy enterprise, takie jak:
SynthID Detector (Google DeepMind), Hive Moderation, Reality Defender oraz Intel
FakeCatcher. Szczególnie interesujące byłoby zbadanie, czy wyższy koszt subskrypcji
przekłada się na rzeczywistą poprawę czułości w stosunku do proponowanego
podejścia wielosygnałowego.

\subsection{Odporność na ataki adwersarialne}

Praktyczny detektor musi być odporny na celowe działania twórców deepfake'ów.
Kierunkiem wartym zbadania jest testowanie detektora wobec technik adversarial,
takich jak dodawanie szumu gaussowskiego przed publikacją, kompresja
H.265 z niskim bitrate oraz preprocessing przez filtry stylistyczne
(np.\ efekt ,,film grain'').
Badania w tym obszarze pozwolą na identyfikację najsłabszych ogniw proponowanego
potoku i ukierunkowanie dalszych prac na ich wzmocnienie.

\subsection{Metody uczenia maszynowego}

Obecna implementacja opiera się na metodach sygnałowych opartych na
regułach heurystycznych. Naturalnym kierunkiem rozwoju jest zastąpienie
lub uzupełnienie modułu decyzyjnego klasyfikatorem uczonym nadzorowanym ---
np.\ lekką siecią CNN na cechach spektralnych lub modelem XGBoost
na wektorze cech numerycznych --- co potencjalnie zwiększy zdolność generalizacji
na nowych modelach generatywnych bez konieczności ręcznej aktualizacji progów.

\subsection{Detekcja w czasie rzeczywistym i integracja z platformami}

Obecna architektura przetwarza pliki wsadowo. Kolejnym krokiem praktycznym
byłoby opracowanie wersji strumieniowej zdolnej do analizy wideo na żywo
(np.\ podczas transmisji na platformach streamingowych) lub wdrożonej jako
wtyczka do CMS redakcyjnych, co odpowiada na rosnące zapotrzebowanie
branży mediowej na narzędzia weryfikacji treści w czasie rzeczywistym.

\subsection{Cyfrowe znaki wodne jako proaktywna ochrona}

Obok reaktywnej detekcji warto rozwinąć badania nad proaktywnym
oznaczaniem treści AI na poziomie modelu generatywnego (C2PA,
SynthID). Analiza skuteczności przetrwania tych znaków po typowych
operacjach edycyjnych (przycięcie, zmiana rozdzielczości, transkodowanie)
stanowi otwarty problem badawczy z bezpośrednim przełożeniem na regulacje prawne UE
dotyczące oznaczania treści syntetycznych~\cite{EUAIAct2024}.
